% --- Archon physics formalization source begin ---
% archon:physics
% archon:covers PhyXMiniProblems/problem_phyx_mini_0682.lean
% archon:source-report reports/phyx_mini/problem_phyx_mini_0682.source.json

\chapter{Physics problem phyx\_mini\_0682}
\label{ch:PhyXMiniProblems_problem_phyx_mini_0682}

\paragraph{Problem source.}
\#\# Physical scenario

A rectangular parallelepiped has dimensions \textbackslash{}( a \textbackslash{}), \textbackslash{}( b \textbackslash{}), and \textbackslash{}( c \textbackslash{}) as shown in figure.

\#\# Question

Notice that \textbackslash{}( \textbackslash{}vec\{R\}\_1 \textbackslash{}), \textbackslash{}( c\textbackslash{}hat\{k\} \textbackslash{}), and \textbackslash{}( \textbackslash{}vec\{R\}\_2 \textbackslash{}) make a right triangle. Obtain a vector expression for the body diagonal vector \textbackslash{}( \textbackslash{}vec\{R\}\_2 \textbackslash{}).

\#\# Answer choices

- A: A: \textbackslash{}( c\textbackslash{}hat\{i\} + b\textbackslash{}hat\{j\} + b\textbackslash{}hat\{k\} \textbackslash{})

- B: B: \textbackslash{}( b\textbackslash{}hat\{i\} + a\textbackslash{}hat\{j\} + c\textbackslash{}hat\{k\} \textbackslash{})

- C: C: \textbackslash{}( a\textbackslash{}hat\{i\} + b\textbackslash{}hat\{j\} + c\textbackslash{}hat\{k\} \textbackslash{})

- D: D: \textbackslash{}( a\textbackslash{}hat\{i\} + c\textbackslash{}hat\{j\} + b\textbackslash{}hat\{k\} \textbackslash{})

\#\# Figure caption (auxiliary; use the image as primary evidence)

The image depicts a 3D rectangular box in a coordinate system with axes labeled as \textbackslash{}(x\textbackslash{}), \textbackslash{}(y\textbackslash{}), and \textbackslash{}(z\textbackslash{}). The box is defined by three dimensions: \textbackslash{}(a\textbackslash{}), \textbackslash{}(b\textbackslash{}), and \textbackslash{}(c\textbackslash{}) along the \textbackslash{}(x\textbackslash{}), \textbackslash{}(y\textbackslash{}), and \textbackslash{}(z\textbackslash{}) axes, respectively. The origin of the coordinate system is marked by \textbackslash{}(O\textbackslash{}). 

Two vectors, \textbackslash{}(\textbackslash{}vec\{R\_1\}\textbackslash{}) and \textbackslash{}(\textbackslash{}vec\{R\_2\}\textbackslash{}), originate from the origin \textbackslash{}(O\textbackslash{}). \textbackslash{}(\textbackslash{}vec\{R\_1\}\textbackslash{}) is directed diagonally towards a corner of the rectangular box. \textbackslash{}(\textbackslash{}vec\{R\_2\}\textbackslash{}) is directed along the diagonal of the base of the rectangular box parallel to the \textbackslash{}(xy\textbackslash{})-plane. A triangular plane is visible within the box, shaded and formed by \textbackslash{}(\textbackslash{}vec\{R\_1\}\textbackslash{}), \textbackslash{}(\textbackslash{}vec\{R\_2\}\textbackslash{}), and the edge along the \textbackslash{}(z\textbackslash{})-axis.

Dashed lines are used to represent hidden edges of the box, indicating depth. Each face of the box is a rectangle, and the entire box is shaded in a light color for visibility.

\#\# Dataset metadata

Kinematics; Spatial Relation Reasoning, Physical Model Grounding Reasoning

\paragraph{Recorded answer/context.}
C: C: \textbackslash{}( a\textbackslash{}hat\{i\} + b\textbackslash{}hat\{j\} + c\textbackslash{}hat\{k\} \textbackslash{})

\paragraph{Figure/image path.}
/root/proposal\_for\_physic/science-mango/phyx\_mini\_run/phyx\_data/test\_image/682.png

\paragraph{Formalization target.}
create a compiling Lean file with sorry bodies at `PhyXMiniProblems/problem\_phyx\_mini\_0682.lean`. The Lean declarations must preserve the physical quantities, dimensions or dimensional roles, figure labels, governing-law hypotheses, and final relation expressed by this problem.
Use Mathlib/Physlib names found through LeanExplore where available. If a physics API is missing, introduce faithful local abstractions rather than scalar placeholder aliases.

\begin{theorem}[Physics formalization target]
\label{thm:physics:phyx_mini_0682:target}
\lean{PhyXMiniProblems.ProblemPhyXMini0682.problem_phyx_mini_0682}
\uses{def:physics:phyx-mini-0682:phyxminiproblems-problemphyxmini0682-lengthreadout, def:physics:phyx-mini-0682:phyxminiproblems-problemphyxmini0682-displacementreadout, def:physics:phyx-mini-0682:phyxminiproblems-problemphyxmini0682-ihat, def:physics:phyx-mini-0682:phyxminiproblems-problemphyxmini0682-jhat, def:physics:phyx-mini-0682:phyxminiproblems-problemphyxmini0682-khat, def:physics:phyx-mini-0682:phyxminiproblems-problemphyxmini0682-rectangularparallelepipedsetup, def:physics:phyx-mini-0682:phyxminiproblems-problemphyxmini0682-matchessuppliedfigure, def:physics:phyx-mini-0682:phyxminiproblems-problemphyxmini0682-haspositiveedgelengths, def:physics:phyx-mini-0682:phyxminiproblems-problemphyxmini0682-satisfiesbasefacediagonalgeometry, def:physics:phyx-mini-0682:phyxminiproblems-problemphyxmini0682-satisfiesbodydiagonalrighttriangle, def:physics:phyx-mini-0682:phyxminiproblems-problemphyxmini0682-matchesdisplayedbodydiagonal}
The assigned autoformalize agent should translate this physics problem into Lean declarations in the covered file, with theorem and lemma proof bodies written as `by sorry`.
\end{theorem}
\begin{proof}
This is an autoformalization task, not a proof task. Produce statements that can later be proved by the physics prover without weakening the physical model.
\end{proof}

% --- Archon declaration topology begin ---
\section{Declaration topology}
\begin{definition}[Length Quantity]
  \label{def:physics:phyx-mini-0682:phyxminiproblems-problemphyxmini0682-lengthquantity}
  \lean{PhyXMiniProblems.ProblemPhyXMini0682.LengthQuantity}
  A nonnegative physical length, independent of the unit used to read it
\end{definition}
\begin{proof}
  This is the declared model definition of Length Quantity.
\end{proof}
\begin{definition}[Spatial Displacement]
  \label{def:physics:phyx-mini-0682:phyxminiproblems-problemphyxmini0682-spatialdisplacement}
  \lean{PhyXMiniProblems.ProblemPhyXMini0682.SpatialDisplacement}
  A dimensionful spatial displacement with three signed Cartesian components
\end{definition}
\begin{proof}
  This is the declared model definition of Spatial Displacement.
\end{proof}
\begin{definition}[length Readout]
  \label{def:physics:phyx-mini-0682:phyxminiproblems-problemphyxmini0682-lengthreadout}
  \lean{PhyXMiniProblems.ProblemPhyXMini0682.lengthReadout}
  \uses{def:physics:phyx-mini-0682:phyxminiproblems-problemphyxmini0682-lengthquantity}
  The scalar readout of a physical length in a coherent unit system
\end{definition}
\begin{proof}
  This is the declared model definition of length Readout.
\end{proof}
\begin{definition}[displacement Readout]
  \label{def:physics:phyx-mini-0682:phyxminiproblems-problemphyxmini0682-displacementreadout}
  \lean{PhyXMiniProblems.ProblemPhyXMini0682.displacementReadout}
  \uses{def:physics:phyx-mini-0682:phyxminiproblems-problemphyxmini0682-spatialdisplacement}
  The Cartesian readout of a physical displacement in a coherent unit system
\end{definition}
\begin{proof}
  This is the declared model definition of displacement Readout.
\end{proof}
\begin{definition}[i Hat]
  \label{def:physics:phyx-mini-0682:phyxminiproblems-problemphyxmini0682-ihat}
  \lean{PhyXMiniProblems.ProblemPhyXMini0682.iHat}
  The dimensionless Cartesian unit vector î along the x axis
\end{definition}
\begin{proof}
  This is the declared model definition of i Hat.
\end{proof}
\begin{definition}[j Hat]
  \label{def:physics:phyx-mini-0682:phyxminiproblems-problemphyxmini0682-jhat}
  \lean{PhyXMiniProblems.ProblemPhyXMini0682.jHat}
  The dimensionless Cartesian unit vector ĵ along the y axis
\end{definition}
\begin{proof}
  This is the declared model definition of j Hat.
\end{proof}
\begin{definition}[k Hat]
  \label{def:physics:phyx-mini-0682:phyxminiproblems-problemphyxmini0682-khat}
  \lean{PhyXMiniProblems.ProblemPhyXMini0682.kHat}
  The dimensionless Cartesian unit vector k̂ along the z axis
\end{definition}
\begin{proof}
  This is the declared model definition of k Hat.
\end{proof}
\begin{definition}[Coordinate Axis]
  \label{def:physics:phyx-mini-0682:phyxminiproblems-problemphyxmini0682-coordinateaxis}
  \lean{PhyXMiniProblems.ProblemPhyXMini0682.CoordinateAxis}
  The three coordinate axes printed in the figure
\end{definition}
\begin{proof}
  This is the declared model definition of Coordinate Axis.
\end{proof}
\begin{definition}[Dimension Label]
  \label{def:physics:phyx-mini-0682:phyxminiproblems-problemphyxmini0682-dimensionlabel}
  \lean{PhyXMiniProblems.ProblemPhyXMini0682.DimensionLabel}
  The three dimension labels printed on the box
\end{definition}
\begin{proof}
  This is the declared model definition of Dimension Label.
\end{proof}
\begin{definition}[Figure Vector Label]
  \label{def:physics:phyx-mini-0682:phyxminiproblems-problemphyxmini0682-figurevectorlabel}
  \lean{PhyXMiniProblems.ProblemPhyXMini0682.FigureVectorLabel}
  The two named vector arrows in the figure
\end{definition}
\begin{proof}
  This is the declared model definition of Figure Vector Label.
\end{proof}
\begin{definition}[Figure Point Label]
  \label{def:physics:phyx-mini-0682:phyxminiproblems-problemphyxmini0682-figurepointlabel}
  \lean{PhyXMiniProblems.ProblemPhyXMini0682.FigurePointLabel}
  Geometrically distinguished points needed to describe the two arrows
\end{definition}
\begin{proof}
  This is the declared model definition of Figure Point Label.
\end{proof}
\begin{definition}[Vector Role]
  \label{def:physics:phyx-mini-0682:phyxminiproblems-problemphyxmini0682-vectorrole}
  \lean{PhyXMiniProblems.ProblemPhyXMini0682.VectorRole}
  The geometric role assigned to each named vector arrow
\end{definition}
\begin{proof}
  This is the declared model definition of Vector Role.
\end{proof}
\begin{definition}[Supplied Parallelepiped Figure]
  \label{def:physics:phyx-mini-0682:phyxminiproblems-problemphyxmini0682-suppliedparallelepipedfigure}
  \lean{PhyXMiniProblems.ProblemPhyXMini0682.SuppliedParallelepipedFigure}
  \uses{def:physics:phyx-mini-0682:phyxminiproblems-problemphyxmini0682-coordinateaxis, def:physics:phyx-mini-0682:phyxminiproblems-problemphyxmini0682-dimensionlabel, def:physics:phyx-mini-0682:phyxminiproblems-problemphyxmini0682-figurevectorlabel, def:physics:phyx-mini-0682:phyxminiproblems-problemphyxmini0682-figurepointlabel, def:physics:phyx-mini-0682:phyxminiproblems-problemphyxmini0682-vectorrole}
  Qualitative labels and marks visible in the primary bitmap. These fields do not impose any component formula on either physical displacement
\end{definition}
\begin{proof}
  This is the declared model definition of Supplied Parallelepiped Figure.
\end{proof}
\begin{definition}[Rectangular Parallelepiped Setup]
  \label{def:physics:phyx-mini-0682:phyxminiproblems-problemphyxmini0682-rectangularparallelepipedsetup}
  \lean{PhyXMiniProblems.ProblemPhyXMini0682.RectangularParallelepipedSetup}
  \uses{def:physics:phyx-mini-0682:phyxminiproblems-problemphyxmini0682-lengthquantity, def:physics:phyx-mini-0682:phyxminiproblems-problemphyxmini0682-spatialdisplacement, def:physics:phyx-mini-0682:phyxminiproblems-problemphyxmini0682-dimensionlabel, def:physics:phyx-mini-0682:phyxminiproblems-problemphyxmini0682-figurevectorlabel, def:physics:phyx-mini-0682:phyxminiproblems-problemphyxmini0682-suppliedparallelepipedfigure}
  The three physical edge lengths and two physical displacement vectors named in the problem, together with the supplied diagram
\end{definition}
\begin{proof}
  This is the declared model definition of Rectangular Parallelepiped Setup.
\end{proof}
\begin{definition}[Matches Supplied Figure]
  \label{def:physics:phyx-mini-0682:phyxminiproblems-problemphyxmini0682-matchessuppliedfigure}
  \lean{PhyXMiniProblems.ProblemPhyXMini0682.MatchesSuppliedFigure}
  \uses{def:physics:phyx-mini-0682:phyxminiproblems-problemphyxmini0682-rectangularparallelepipedsetup}
  The axis assignments and arrow roles read from the primary figure. In particular, this identifies R₁ as the base diagonal and R₂ as the body diagonal, without assuming the requested component expression for R₂
\end{definition}
\begin{proof}
  This is the declared model definition of Matches Supplied Figure.
\end{proof}
\begin{definition}[Has Positive Edge Lengths]
  \label{def:physics:phyx-mini-0682:phyxminiproblems-problemphyxmini0682-haspositiveedgelengths}
  \lean{PhyXMiniProblems.ProblemPhyXMini0682.HasPositiveEdgeLengths}
  \uses{def:physics:phyx-mini-0682:phyxminiproblems-problemphyxmini0682-lengthreadout, def:physics:phyx-mini-0682:phyxminiproblems-problemphyxmini0682-dimensionlabel, def:physics:phyx-mini-0682:phyxminiproblems-problemphyxmini0682-rectangularparallelepipedsetup}
  The physical branch in which all three edge dimensions are nonzero
\end{definition}
\begin{proof}
  This is the declared model definition of Has Positive Edge Lengths.
\end{proof}
\begin{definition}[Satisfies Base Face Diagonal Geometry]
  \label{def:physics:phyx-mini-0682:phyxminiproblems-problemphyxmini0682-satisfiesbasefacediagonalgeometry}
  \lean{PhyXMiniProblems.ProblemPhyXMini0682.SatisfiesBaseFaceDiagonalGeometry}
  \uses{def:physics:phyx-mini-0682:phyxminiproblems-problemphyxmini0682-lengthreadout, def:physics:phyx-mini-0682:phyxminiproblems-problemphyxmini0682-displacementreadout, def:physics:phyx-mini-0682:phyxminiproblems-problemphyxmini0682-ihat, def:physics:phyx-mini-0682:phyxminiproblems-problemphyxmini0682-jhat, def:physics:phyx-mini-0682:phyxminiproblems-problemphyxmini0682-rectangularparallelepipedsetup}
  Vector addition across the rectangular xy base face: its diagonal R₁ is the sum of the a edge along î and the b edge along ĵ
\end{definition}
\begin{proof}
  This is the declared model definition of Satisfies Base Face Diagonal Geometry.
\end{proof}
\begin{definition}[Satisfies Body Diagonal Right Triangle]
  \label{def:physics:phyx-mini-0682:phyxminiproblems-problemphyxmini0682-satisfiesbodydiagonalrighttriangle}
  \lean{PhyXMiniProblems.ProblemPhyXMini0682.SatisfiesBodyDiagonalRightTriangle}
  \uses{def:physics:phyx-mini-0682:phyxminiproblems-problemphyxmini0682-lengthreadout, def:physics:phyx-mini-0682:phyxminiproblems-problemphyxmini0682-displacementreadout, def:physics:phyx-mini-0682:phyxminiproblems-problemphyxmini0682-khat, def:physics:phyx-mini-0682:phyxminiproblems-problemphyxmini0682-rectangularparallelepipedsetup}
  The stated right-triangle law. Traversing R₁ and then the vertical edge c k̂ gives R₂; the base diagonal is perpendicular to that vertical leg. Neither field contains the requested three-component expansion of R₂
\end{definition}
\begin{proof}
  This is the declared model definition of Satisfies Body Diagonal Right Triangle.
\end{proof}
\begin{definition}[Answer Choice]
  \label{def:physics:phyx-mini-0682:phyxminiproblems-problemphyxmini0682-answerchoice}
  \lean{PhyXMiniProblems.ProblemPhyXMini0682.AnswerChoice}
  Labels of the four displayed multiple-choice expressions
\end{definition}
\begin{proof}
  This is the declared model definition of Answer Choice.
\end{proof}
\begin{definition}[displayed Body Diagonal Readout]
  \label{def:physics:phyx-mini-0682:phyxminiproblems-problemphyxmini0682-displayedbodydiagonalreadout}
  \lean{PhyXMiniProblems.ProblemPhyXMini0682.displayedBodyDiagonalReadout}
  \uses{def:physics:phyx-mini-0682:phyxminiproblems-problemphyxmini0682-lengthreadout, def:physics:phyx-mini-0682:phyxminiproblems-problemphyxmini0682-ihat, def:physics:phyx-mini-0682:phyxminiproblems-problemphyxmini0682-jhat, def:physics:phyx-mini-0682:phyxminiproblems-problemphyxmini0682-khat, def:physics:phyx-mini-0682:phyxminiproblems-problemphyxmini0682-rectangularparallelepipedsetup, def:physics:phyx-mini-0682:phyxminiproblems-problemphyxmini0682-answerchoice}
  The vector expression printed beside each answer choice
\end{definition}
\begin{proof}
  This is the declared model definition of displayed Body Diagonal Readout.
\end{proof}
\begin{definition}[recorded Dataset Answer]
  \label{def:physics:phyx-mini-0682:phyxminiproblems-problemphyxmini0682-recordeddatasetanswer}
  \lean{PhyXMiniProblems.ProblemPhyXMini0682.recordedDatasetAnswer}
  \uses{def:physics:phyx-mini-0682:phyxminiproblems-problemphyxmini0682-answerchoice}
  The answer label recorded in the dataset metadata
\end{definition}
\begin{proof}
  This is the declared model definition of recorded Dataset Answer.
\end{proof}
\begin{definition}[Matches Displayed Body Diagonal]
  \label{def:physics:phyx-mini-0682:phyxminiproblems-problemphyxmini0682-matchesdisplayedbodydiagonal}
  \lean{PhyXMiniProblems.ProblemPhyXMini0682.MatchesDisplayedBodyDiagonal}
  \uses{def:physics:phyx-mini-0682:phyxminiproblems-problemphyxmini0682-displacementreadout, def:physics:phyx-mini-0682:phyxminiproblems-problemphyxmini0682-rectangularparallelepipedsetup, def:physics:phyx-mini-0682:phyxminiproblems-problemphyxmini0682-answerchoice, def:physics:phyx-mini-0682:phyxminiproblems-problemphyxmini0682-displayedbodydiagonalreadout}
  A displayed expression agrees with the physical body diagonal in all units
\end{definition}
\begin{proof}
  This is the declared model definition of Matches Displayed Body Diagonal.
\end{proof}
% --- Archon declaration topology end ---
% --- Archon physics formalization source end ---
